% TOP_PROBLEM: {"id": 1101002, "title": "Questions in Geometric Group Theory — Q 10.2", "category_id": 17, "category": {"id": 17, "name": "group_theory", "display_name": "Group Theory", "description": "Problems about groups, group actions, representations, and related algebraic structures.", "slug": "group-theory", "order_index": 17, "created_at": "2026-07-31T15:26:25.671Z"}, "status": "open", "problem_number": "AMR-010-1002", "set_id": 13, "statusQualification": "Retains the source general CW-complex domain, without adding finiteness or asphericity."}
% FIRST_POSTED: 2026-10-06
% ATTEMPT_STATUS: unresolved
% OPENAI_CLAIM: {"schema_version":1,"source_repo":"https://github.com/openai/math","source_commit":"adc7f1241b42e322a6451854ab7e4b4c146bf78a","release_date":"2026-10-06","match":"related","resolution":"partial","independently_reviewed":false,"second_pass":"n/a","families":[{"family":"321","title":"A counterexample to Wall's finite D(2) problem","subject":"Topology","manuscripts":[{"title":"A Counterexample to Wall's D(2) Problem","pdf_path":"preprints/A-Counterexample-to-Walls-D2-Problem-October-6-2026/wall-d2-counterexample.pdf","readme_path":"preprints/A-Counterexample-to-Walls-D2-Problem-October-6-2026/README.md","date":"2026-10-06","theorem_ref":"thm:main (sections/01-introduction.tex:19)","statement_sources":[{"path":"preprints/A-Counterexample-to-Walls-D2-Problem-October-6-2026/build/sections/01-introduction.tex","line":19,"label":"thm:main"}]}],"lean":null,"reasoning_trace":null}]}
\documentclass[11pt]{article}
\usepackage[margin=25mm]{geometry}
\usepackage{fontspec}
\setmainfont{Latin Modern Roman}
\usepackage{amsmath,amssymb,mathtools}
\usepackage{unicode-math}
\setmathfont{Latin Modern Math}
\usepackage{enumitem}
\usepackage{xurl,hyperref,bookmark}
\hypersetup{colorlinks=true,urlcolor=blue}
\setcounter{secnumdepth}{0}
\setlength{\parindent}{0pt}
\setlength{\parskip}{5pt}
\setlength{\emergencystretch}{3em}
\begin{document}
\section*{Questions in Geometric Group Theory — Q 10.2}

\subsection{OpenAI's claim}
Wall's finite D(2) question is claimed to have a negative example: a finite three-dimensional complex with infinite fundamental group satisfies D(2) but has no finite two-dimensional homotopy model. The cited manuscript argues for that result; its relation to this problem is specified below.

\subsection{Mathematical statement of the claimed result}
\paragraph{Mathematical claim 1.} OpenAI claims a finite connected CW complex $X$ of dimension three such that, for its universal cover $\widetilde X$,
\[
H_i(\widetilde X;\mathbb Z)=0\quad(i>2),\qquad
H^3(X;M)=0\quad\text{for every }\mathbb Z[\pi_1X]\text{-module }M,
\]
yet $X$ is without a homotopy equivalent finite CW model in dimension two or below. The cohomology condition ranges over all local coefficient modules, and no asphericity assumption is made.

\subsection{Outline of the argument}
Manuscript 1 (family 321, 2026-10-06): A rank-one twisted-homology obstruction is proved by repeatedly lifting a fixed finite source through cyclic-cover image towers, whose simplex growth cannot continue. A five-generator presentation has a finitely normally generated perfect subgroup and a quotient involution detected by a complex character. Its explicit boundary matrix gives vanishing twisted second homology; finite plus construction preserves this while contradicting the necessary two-dimensional character condition.

\subsection{Scope relative to this problem}
The manuscript constructs a finite three-complex satisfying the local-coefficient cohomology condition and excludes every finite two-complex homotopy model. This page allows a two-complex with infinitely many cells, whose nonexistence is not claimed by the paper. Consequently the constructed example is not asserted to meet the target's stronger nonexistence requirement. It is a related result on the finite-target formulation.

\subsection{Formal verification}
Family 321 has no Lean scope page in this snapshot; no formalization of these manuscript results is listed.

This site has not executed the Lean files or checked their dependencies. These descriptions report the scope of the linked declarations, rather than a new verification of the proofs.

\subsection{Status caveat}
OpenAI attributes these results to an internal model. Its release notice cautions that results lacking full formalization can contain mistakes. This site has not independently assessed the proofs; a solved label records the scope of the reported claim.

\subsection{Sources}
\begin{itemize}
\item OpenAI, \emph{A Counterexample to Wall's D(2) Problem}, 2026-10-06; citation key \texttt{OAI:A-Counterexample-to-Walls-D2-Problem-October-6-2026}. \href{https://github.com/openai/math/blob/adc7f1241b42e322a6451854ab7e4b4c146bf78a/preprints/A-Counterexample-to-Walls-D2-Problem-October-6-2026/wall-d2-counterexample.pdf}{PDF (GitHub)}; \href{https://github.com/openai/math/raw/adc7f1241b42e322a6451854ab7e4b4c146bf78a/preprints/A-Counterexample-to-Walls-D2-Problem-October-6-2026/wall-d2-counterexample.pdf}{PDF (direct)}; \href{https://github.com/openai/math/blob/adc7f1241b42e322a6451854ab7e4b4c146bf78a/preprints/A-Counterexample-to-Walls-D2-Problem-October-6-2026/README.md}{Manuscript page and citation}; \href{https://github.com/openai/math/blob/main/preprints/A-Counterexample-to-Walls-D2-Problem-October-6-2026/wall-d2-counterexample.pdf}{Latest version}; Statement source: \href{https://github.com/openai/math/blob/adc7f1241b42e322a6451854ab7e4b4c146bf78a/preprints/A-Counterexample-to-Walls-D2-Problem-October-6-2026/build/sections/01-introduction.tex\#L19}{thm:main}.
\item \href{https://github.com/openai/math/blob/adc7f1241b42e322a6451854ab7e4b4c146bf78a/CONTENTS.md}{OpenAI catalogue entry}.
\end{itemize}
\end{document}
